%In this study, we aim to address the challenge of learning from a dataset with noisy labels, focusing specifically on the context of graph classification. Here, we outline the various areas related to our research.
\textbf{Learning under label noise.} A large body of work is devoted to the challenge of learning with noisy labels. Several methods are based on \textbf{robust loss functions}, using symmetric losses \citet{ghosh2017robust}, or loss correction methods \citep{patrini2017making}. Other works are the \textbf{Neighboring-based noise identification} approaches \citep{zhu2022detecting}.
%\cite{zhu2021clusterability} and \cite{zhu2022detecting} in particular propose methods based on the clusterability condition. 
Several approaches are based on the \textbf{early learning phenomenon} \citet{arpit2017closer}, and others proposed to improve the quality of training data by treating samples with a small loss value as correctly labeled during the training process \citep{gui2021towards}. 
Additional methods for learning in the presence of noisy labels are in Appendix~\ref{sec:additional_related}.

\textbf{Graph Learning in noisy scenarios.} 
The methods discussed earlier focus mainly on learning from noisy labels in image datasets. Unlike images, graphs exhibit noise in labels, graph topology (e.g., adding/removing edges or nodes), and node features. Most prior work discusses noisy node labels \citet{graphC1, contrastivelearning5, omg7, ssl8, cl9, nrgnn11,Kang2018RobustGL}, while noise at the edge and feature levels have also been explored \citep{structuralnoise2, structuralnoise3, ssl8}.
However, fewer studies investigate graph classification under noise, limiting progress in applying and improving graph classification tasks. The seminal work of \citet{graphC1} addresses graph classification with label noise, proposing a surrogate loss to discard noisy labels under certain assumptions, without comparison to clean label scenarios. More recently, \citet{omg7} introduces a method combining contrastive learning and MixUp \citet{mixup2} within the loss function to improve generalization, along with a curriculum learning strategy to dynamically discard noisy samples.
In contrast, we propose tackling noisy labels using an effective loss function inspired by \citet{wani2023combining} or by enhancing graph smoothness.


% The methods described in the previous section primarily address learning from noisy labels in image datasets. Unlike other data types, graphs exhibit noise in various forms: within labels, graph topology modifications (such as adding or removing edges or nodes), and node features. 
% Most of the works presented in the literature focus on noisy labels for the nodes \cite{graphC1, pairwiseinteraction4, contrastivelearning5, omg7, ssl8, cl9, unionnet10, nrgnn11}. Noise present at the edge level and in features has also been studied in the literature  \cite{structuralnoise2, structuralnoise3, ssl8}.  
% In this scenario, fewer studies have explored graph classification (\cite{graphC1, ssl8}). This disparity limits the advancement of new algorithms that could enhance graph classification, thereby restricting the potential benefits for the broader community.  A seminal work on the topic, \cite{graphC1}, concerns graph classification under label noise, proposing a surrogate loss that aims at discarding labels that are estimated to be noisy. Their approach works only under certain assumptions, and no comparisons are made with clean-label scenarios. More recently in \cite{omg7}, a novel method that concerns graph classification under label noise has been studied. Their approach relies on contrastive learning like most of the node classification methods we reviewed, they use a MixUp \cite{mixup1, mixup2} approach within the loss function. The positive, and negative graph representation pairs are mixed to learn a representation that tries to improve their generalization power. %Moreover, they include an additional loss term that has the objective of learning to discriminate those samples that are probably affected by noise. In order to do so they include pseudo-labels biased by the neighboring graphs of each sample. Through this optimization aspect it is possible to discard some sample from the training set as the learning phase proceeds. 
% They also use a curriculum learning strategy to select dynamically the number of expected noisy samples to discard.
% In contrast to these approaches, we address noise by designing an effective loss function inspired by \cite{wani2023combining} and also propose enhancing graph smoothness by removing negative eigenvalues from the weight matrix.\\
 %Our approach is since it is a noise-robust and task-agnostic loss function. We augment its original design with some slight differences that enhance its performance in graph classification.\\
% Dirichlet Energy literature.
\textbf{Dirichlet Energy, Smoothing bias and Sharpening} 
%Graph Neural Networks (GNNs) face several challenges, including limited message-passing expressiveness \citep{WL2}, over-smoothing \citep{oversmoothing2}, and over-squashing \citep{oversquashing3}. 
Dirichlet energy is a key measure in GNN, quantifying the smoothness or variation of features across nodes \citep{dirichlet}. 
%Over-smoothing has been studied using Dirichlet energy \citep{dirichlet}, which quantifies signal smoothness across graph nodes. 
Most GNNs function as low pass filters, emphasizing low frequency components while diminishing high frequency ones \citep{nt2019revisiting, rusch2023survey}. Specifically, \citet{nt2019revisiting} showed this phenomenon holds for graphs without non trivial bipartite components, with self loops further shrinking eigenvalues. \citet{noteonoversmoothing, oversmoothing2} prove that GNN Dirichlet energy exponentially decreases with additional layers when the product of the largest singular value of the weight matrix and the largest eigenvalue of the normalized Laplacian is less than one. GNN learnable weight matrices fundamentally control whether features are smoothed or sharpened \citep{GRAFF}. 
While Dirichlet energy evolution has been studied in relation to oversmoothing \citet{noteonoversmoothing, nt2019revisiting}, and various mitigation approaches leveraging energy properties exist \citet{FAGCN, dirichlet_constraint, EEConv}, these work has primarily focused on node classification, where oversmoothing significantly impacts performance \citep{2sides}. The role of energy dynamics in graph classification, especially with label noise, remains less explored.
In this work, we provide theoretical and practical insights on leveraging Dirichlet energy to enhance graph classification performance, even in the presence of label noise (comprehensive discussion in Appendix \ref{sec:additional_related}).


% \textbf{Smoothing bias} 
% Most GNNs function as low-pass filters, emphasizing low-frequency components while diminishing high-frequency ones \citep{nt2019revisiting, rusch2023survey}. 
%Specifically, \cite{nt2019revisiting} showed this phenomenon holds for graphs without non-trivial bipartite components, with self-loops further shrinking eigenvalues. 
%\cite{noteonoversmoothing, oversmoothing2} prove that GNN Dirichlet energy exponentially decreases with additional GCN layers when the product of the largest singular value of the weight matrix and the largest eigenvalue of the normalized Laplacian is less than one. 
%Similarly, \cite{oversquashing1} finds that non-bipartite graphs, especially without residual connections, exhibit low-frequency dominance. They also show that continuous-time models like CGNN, GRAND, and PDE-GCND maintain low-pass filtering. \cite{noteonoversmoothing, oversmoothing2} prove that GNN Dirichlet energy exponentially decreases with additional GCN layers when the product of the largest singular value of the weight matrix and the largest eigenvalue of the normalized Laplacian is less than one. Here, it is important to emphasize that \cite{Kang2018RobustGL} examines graph classification under label noise using the mix-up technique. While the mix-up may indirectly promote smoothness in the graph, they do not discuss or establish a relationship between graph smoothness and the Dirichlet energy. Furthermore, their work centers on the smoothness of clusters within the graphs, rather than on the overall smoothness of the graph structure. 

% A more comprehensive discussion of related works is given in Appendix \ref{sec:additional_related}.

% Most GNNs act as low-pass filters, highlighting low-frequency components in the feature space while suppressing high-frequency ones. This behavior, extensively studied in the literature (e.g., \citet{nt2019revisiting, wu2019simplifying, rusch2023survey}), is shown by \citet{nt2019revisiting} to hold for all graphs without non-trivial bipartite components. The low-pass effect is further enhanced by adding self-loops, as they shrink the eigenvalues toward zero. 
% Similarly, \citet{oversquashing1} demonstrate that non-bipartite graphs, especially those lacking residual connections, display low-frequency dominant dynamics across various initial conditions. They also show that certain continuous-time message-passing networks, like CGNN, GRAND, and PDE-GCND, retain low-pass filtering properties, preventing high-frequency dominance.

% \citet{noteonoversmoothing, oversmoothing2} theoretically analyzes the over-smoothing phenomenon in GNNs, showing that under certain conditions on the weight matrix, the Dirichlet energy decreases exponentially as more GCN layers are added. This happens when the product of the largest singular value of the weight matrix and the largest eigenvalue of the normalized Laplacian is less than one.
%\cite{noteonoversmoothing, oversmoothing2} also studied over-smoothing phenomena for GNNs from a theoretical point of view. These papers prove that under specific conditions on the weight matrix, the Dirichlet energy decreases exponentially with additional GCN layers. This occurs when the largest singular value of the weight matrix, multiplied by the largest eigenvalue of the normalized Laplacian, is less than one.

 % A more comprehensive discussion of related work is relegated to Section \ref{sec:additional_related} in the Appendix.
% Related work is discussed in detail in Appendix Section \ref{sec:additional_related}.